\section*{Open Questions}

The following five questions remain genuinely open after this replication.
Each is grounded in a specific gap identified during the paper re-read; none
is a rhetorical prompt.

\subsection*{Q1 --- MMEJ upregulation at high LET}

\textbf{Question.} Does the model predict alt-EJ / MMEJ upregulation at high LET,
and if not, how badly does the missing MMEJ branch bias the residual-foci
curves at 100+~keV/$\mu$m?

\textbf{Basis.} The paper's Discussion explicitly lists alt-EJ / MMEJ as an
omitted pathway. Empirical literature (Sharma et al. 2015, Truong et al. 2013)
shows MMEJ can carry $>$30\% of end-joining at high LET where classical NHEJ
is overwhelmed. Belov's Fig 10 (high-LET $\gamma$-H2AX kinetics) attributes all
slow-clearance behaviour to reduced NHEJ efficiency, but the same signature is
predicted by a missing MMEJ compartment; the two are not identifiable from a
$\gamma$-H2AX read-out alone.

\textbf{Next steps.} Add a 4th ODE branch with MMEJ-specific PARP1/POLQ kinetics
(rate constants from Sallmyr \& Tomkinson 2018). Refit against Grabarz et al.\
2013 MMEJ-reporter data (I-SceI + POLQ knockdown). Compare LET-dependence of
MMEJ fraction predicted by the 4-pathway model vs the paper's 3-pathway fit;
report $\Delta$AIC and identifiability of MMEJ rate constants under the
$\gamma$-H2AX-only observation model.

\subsection*{Q2 --- Resection-checkpoint identifiability}

\textbf{Question.} Are the resection-checkpoint rate constants (K5, K6 governing
MRN/CtIP-mediated 5$'$ end resection) practically identifiable from the datasets
Belov fits, or are they in a flat likelihood valley with the HR/SSA branching
ratios?

\textbf{Basis.} The paper fits 46 rate constants against a modest number of
literature time-courses using Newton--Raphson (\S3). No profile-likelihood,
Fisher information, or Sobol sensitivity analysis is reported. With 46
parameters and $\sim$10 time-courses, practical unidentifiability is nearly
certain. The resection step is the HR/SSA branch point; if K5, K6 are
non-identifiable then HR-vs-SSA competition predictions have no error bars.

\textbf{Next steps.} Run PottersWheel-style profile-likelihood on the fitted
parameter vector using the 12 literature time-courses in \S3 as the training set.
Report 95\% CI on each rate constant. For resection-checkpoint constants
specifically, fit against DSB-end-resection assays (Zhou et al.\ 2014, RPA-ssDNA
IF) to pin K5, K6 independently, then propagate to the $\gamma$-H2AX predictions
and quantify the reduction in posterior spread on the Fig 11 ratios.

\subsection*{Q3 --- Extension to cancer cells with pathway defects}

\textbf{Question.} Can the model be extended to cancer cell lines with defined
pathway defects (BRCA1-null, RAD51 paralogue mutants, POLQ-high) without re-fitting
all 46 constants, or does each defect require a full re-fit?

\textbf{Basis.} The paper handles four repair-deficient classes (DNA-PKcs$^-$,
LigIV$^-$, BRCA2$^-$, ERCC1/XPF$^-$) as parameter shifts on the same 46-constant
vector (Table A.1 footnotes describe which rate constants are zeroed for each).
Whether that ansatz generalises to cancer cell lines with more subtle,
dose-dependent, or synthetic-lethal pathway defects is untested. Clinically
relevant BRCA1-null and PARP-inhibited backgrounds are absent.

\textbf{Next steps.} Take three cancer lines with well-characterised repair
defects (MDA-MB-436 BRCA1-null; PEO1 BRCA2-null; HCT116 POLB-null) and their
isogenic revertants. Fit only the defect-associated rate-constant multipliers
(leave all other 46 constants fixed at Belov's values). Report goodness-of-fit
vs a full 46-parameter re-fit; if the parameter-shift ansatz holds to within
1$\sigma$ of foci counts at 2 / 6 / 24~h, the model has clinical extrapolation
value; if not, the model is fit-restricted to the training cell lines.

\subsection*{Q4 --- In-vivo vs in-vitro rate discrepancy for Ku binding}

\textbf{Question.} How large is the in-vivo vs in-vitro rate discrepancy for
the K1..K7 Ku-binding constants, and does the units-typo hypothesis close the
gap or open a new one?

\textbf{Basis.} Table A.1 K1..K7 as printed imply Ku$\to$DSB half-time
$\approx 4.6\times 10^6$~min; in-vitro FRAP data (Reynolds et al.\ 2012,
Mari et al.\ 2006) report $\sim 15$--$30$~s in living cells. A 6--7 order-of-magnitude
gap. If the correct units multiplier is $10^6$ (i.e.\ $\mu$M$^{-1}$ vs M$^{-1}$
\emph{and} min$^{-1}$ vs s$^{-1}$ combined), the model becomes internally
consistent; but nobody has empirically verified this on in-vivo data.

\textbf{Next steps.} Digitise the Ku80-GFP FRAP recovery curves from Mari 2006
(Fig 2C) and Reynolds 2012 (Fig 4B) using WebPlotDigitizer. Fit K1 as a free
parameter against those curves under two hypotheses: (H1) units-as-printed,
(H2) K1 rescaled by $10^6$. Report the reduced $\chi^2$ for each; publish
whichever units convention is consistent with the FRAP data as a formal erratum
to Table A.1. Independently contact corresponding author dem@jinr.ru for the
driver code (only asked-for artefact that would close this in one email).

\subsection*{Q5 --- Coupling to chromosome-aberration models}

\textbf{Question.} Can the $\gamma$-H2AX foci output be coupled to a
chromosome-aberration model (dicentrics, translocations) to predict downstream
cytogenetic damage, or is $x_{14}$ a dead-end read-out?

\textbf{Basis.} Belov's model stops at $\gamma$-H2AX foci; there is no coupling
to mis-rejoining, chromosomal exchange, or aberration formation. Yet the LUCID
initiative's downstream use case is exactly that: predicting late cytogenetic
and carcinogenic endpoints from early DSB kinetics. The BIANCA model
(Ballarini 2014), MEDRAS (McMahon 2016), and NASA's NSCR-2020 all include an
aberration-formation step that Belov's framework omits.

\textbf{Next steps.} Extend the model with a Ballarini-style pairwise-misrejoining
rule operating on the $(x_{ij}, y_{ij}, z_{ij})$ intermediate pools. Predict
dicentric-per-cell as a function of dose and LET; validate against Loucas \&
Cornforth 2001 (Fe-56 dicentrics at LET 200~keV/$\mu$m) and Cornforth 2013
(proton-vs-$\gamma$ dicentric RBE). Report whether Belov's fitted rate constants
produce a plausible RBE-LET curve for dicentrics, or whether the model must be
re-fit against cytogenetic endpoints as well.
